\documentclass[11pt]{article}
\usepackage[utf8]{inputenc}
\usepackage[T1]{fontenc}
\usepackage{amsmath,amssymb,amsthm}
\usepackage{geometry}
\geometry{margin=1in}
\usepackage{hyperref}
\hypersetup{colorlinks=true,linkcolor=blue,urlcolor=blue}

\newtheorem{thm}{Theorem}[section]
\newtheorem{defn}[thm]{Definition}
\newtheorem{prop}[thm]{Proposition}
\newtheorem{cor}[thm]{Corollary}
\theoremstyle{remark}
\newtheorem{rem}[thm]{Remark}

\newcommand{\OpSet}{\textsf{OpSet}}
\newcommand{\eq}{\approx}
\newcommand{\N}{\mathbb{N}}
\newcommand{\R}{\mathbb{R}}
\newcommand{\Z}{\mathbb{Z}}
\newcommand{\Q}{\mathbb{Q}}
\newcommand{\C}{\mathbb{C}}

\title{\textbf{VR-SetsOp.}\\[2pt]
\large An Operational Set Theory Unifying ZFC and ZFA (Brouwer Edition)\\[2pt]
\normalsize \emph{(formerly titled VR-Sets; same Zenodo concept 10.5281/zenodo.20303535 --- retitled to name the operational carrier \OpSet{})}}
\author{Vitaly Reznik}
\date{2026 --- Version 1.2.0}

\begin{document}
\maketitle

\begin{abstract}
We reconstruct set theory inside the VR programme as an operational universe: a set is a
\emph{functionality} that, on a query, reveals its members --- formally, a pointed graph --- and
the identity of sets is a \emph{witnessed bisimulation}, carried as a relation and never collapsed
to a quotient. On this carrier we obtain, machine-checked in Lean 4 and almost entirely below the
\texttt{Classical.choice} floor: membership, strong extensionality (identity is exactly
co-membership), the empty set, pairing, union, the power set, the von Neumann naturals and $\omega$,
and the Separation and Replacement schemas --- every axiom of ZF without Foundation. The power set is
machine-checked \emph{choice-free} but is \emph{impredicative}, which isolates the obstruction to
power set in the operational register as impredicativity, not choice. Anti-Foundation is a theorem ---
indeed axiom-free: a member is the same graph re-pointed, so nesting (including self-membership) is
native, and decoration of any graph is almost definitional. Well-foundedness is not an axiom
but a predicate
\texttt{IsGrounded}: the ZFC universe and the anti-foundational (Aczel) universe are two fragments
of one operational domain, not a foundational choice one is forced to make. Version~1.0.3 corrected
the central error of earlier versions: the earlier text claimed the operational universe is countable
and inferred that the Axiom of Choice is a free theorem and that paradoxes such as Banach--Tarski
cannot arise --- a Bishop reading that contradicts the operational continuum the cycle built on the
Brouwer path. We replace it: the operational universe --- sets as becoming --- is genuinely
non-enumerable (a machine-checked, choice-free diagonal), while the describable register (sets given
by a finite rule) is countable. The Axiom of Choice is therefore not an operational theorem; only
dependent choice is. The absence of Banach--Tarski is re-explained through conservativity, not
countability.
\\[3pt]
\emph{Keywords:} operational set theory, non-well-founded sets, bisimulation, strong
extensionality, Brouwer, choice sequences, axiom of choice, constructive foundations, Lean 4.
\\[3pt]
\textbf{Version note (1.2.0).} No mathematics is altered, added, or removed. Two changes. (i) The
work is retitled \emph{VR-SetsOp} to name its operational carrier \OpSet{} (the earlier title
\emph{VR-Sets} named the domain; the content, since Version~1.0.3, is entirely the operational
universe); the Zenodo concept (10.5281/zenodo.20303535) is unchanged. (ii) A positioning paragraph is
added to \S\ref{sec:position} stating the exact relation to ZFA: with VR's formal register on
(classical logic, choice) the set layer \emph{realises} Aczel's ZFA and is \emph{equal to it in
strength} (equiconsistent, at the level of ZF --- neither stronger nor weaker), with ZFC recovered as
the well-founded inner class \texttt{IsGrounded}; without the formal register the operational core is
strictly weaker (constructive, choice-free). This resolves the imprecision in the Version~1.1.0 phrase
``holds ZFC and ZFA natively'' and records the one strict advantage at equal strength --- nesting is a
theorem in VR where ZFA posits it.

\textbf{Version note (1.1.0).} A minor version: new mathematics is added, none altered or removed.
Adds the operational \textbf{power set} (\S3.3): $\wp\,x$ with
$z\in\wp\,x\iff z\subseteq x$, machine-checked \emph{choice-free} (\texttt{propext}). With it,
\OpSet{} verifies every axiom of ZF without Foundation. The construction is below the choice floor but
\emph{impredicative} --- this isolates the obstruction to power set in the operational register as
impredicativity, not choice, and corrects the blanket ``no power-set operation'' of the earlier \S7
(Remark~3.10). It also states \textbf{Anti-Foundation as a theorem} on \OpSet{} (\S4.1),
fully axiom-free, with the cost-of-quotient bridge to VR-Sets-ZFA; anchors \S6.4 in Solovay's model;
and names two prices (the setoid discipline; impredicativity's invisibility to the axiom audit, P-1).
No prior mathematics is altered.
\\[3pt]
\textbf{Version note (1.0.4).} Editorial synchronisation, no mathematics altered: (i) Remark~5.2
is corrected --- VR-Numbers' $\R_{VR}$ is the \emph{computable, countable} reals (its Lean
isomorphism with the full classical line is an inexpressibility artifact, VR-Numbers \S VIII.6), the
genuinely non-enumerable object is the operational becoming-continuum (and this universe, \S6.1),
and classical $\R$ is a third thing; (ii) references carry correct, separated DOIs (the formal-system
\emph{paper} is DOI 10.5281/zenodo.20633314, concept 10.5281/zenodo.20212091; 20324240 is its
Lean companion); (iii) the describable register is noted to be a minimal representative of ``finite
rule'' (\S6.2); (iv) the single mathlib borrowing is flagged already in \S1.
\end{abstract}

\tableofcontents

%======================================================================
\section{Introduction and the change since Version 1.0.2}

VR-Sets realises the set-theoretic layer of the VR programme. A set is not a ``thing that contains
elements''; it is an \emph{operation} --- a functionality that, when queried, reveals members, which
are themselves functionalities. This reading is constitutive, not decorative: every notion below
(membership, identity, the operations, well-foundedness) is defined in its terms, and every claim is
checked by the Lean~4 kernel with its axiom cost recorded.

Versions through 1.0.2 presented the operational universe through a classical surrogate (mathlib's
well-founded \texttt{ZFSet}) and argued, in their \S VI, that the universe is countable, whence the
Axiom of Choice holds as a theorem and the classical ``pathologies'' (non-measurable sets,
Banach--Tarski) do not arise. Two things were wrong. First, countability is false of the surrogate
(\texttt{ZFSet} is the full classical uncountable universe), so the claim was carried by prose, not by
the machine. Second, and more deeply, the countability reading is a \emph{Bishop} reading, and it
is incompatible with the operational continuum the VR cycle actually constructed, which took the
\emph{Brouwer} path: there the space of ``becoming'' (lawless choice sequences) is non-enumerable,
choice-free, and Cantor's diagonal works.

We resolve the tension on the side the cycle already chose. The operational universe of sets is
rebuilt from scratch, as functionalities with witnessed-bisimulation identity (\S2), and we prove
(\S6) that this universe --- sets as becoming --- is not enumerable, by a self-contained diagonal
that is choice-free. The countable register is the register of \emph{descriptions}: sets given by a
finite rule. The contrast \emph{done countable / becoming non-enumerable} is the corrected \S VI.
The Axiom of Choice loses its status as a free theorem; this is not a loss but the honest price of a
genuinely uncountable operational continuum, and it aligns VR-Sets with VR-Forms (conservativity)
and with the operational-continuum work.

\paragraph{What is machine-checked.} The development lives in the Lean~4 formalisation of the VR
cycle, module \texttt{VRCycle.SetsOp} (files \texttt{Pointed}, \texttt{Builder}, \texttt{Closure},
\texttt{Omega}, \texttt{Extensionality}, \texttt{Congruence}, \texttt{Becoming},
\texttt{Describable}, \texttt{Schemas}, \texttt{Grounded}, \texttt{Power}, \texttt{AFA}). It is self-contained: it imports
nothing from mathlib for its core --- precisely so that its axiom profile is visible --- the
\emph{single} mathlib borrowing (the pairing inversion in the describable register) being flagged
in Remark~6.5. Throughout, we name the Lean object backing each statement. Almost every object is
axiom-free or depends only on \texttt{propext}; the single exception is flagged honestly in \S6.

%======================================================================
\section{The carrier: sets as revealing functionalities}

\subsection{Operational sets}
\begin{defn}[Operational set; \texttt{OpSet}]
An operational set is a pointed graph $x=(V,E,\mathrm{pt})$ with $V:\mathrm{Type}$,
$E:V\to V\to\mathrm{Prop}$, $\mathrm{pt}:V$. Read $E\,a\,b$ as ``querying the functionality at
vertex $b$ reveals $a$ as a member''; $\mathrm{pt}$ is the set itself. The member of $x$ sitting at
a vertex $a$ is the same graph re-pointed there, $x.\mathrm{child}\,a:=(V,E,a)$. This is exactly the
language in which Aczel's anti-foundation axiom is stated (decoration of a graph), but we take
\emph{no quotient}: the identity of sets is supplied operationally, below.
\end{defn}

\subsection{Operational identity: witnessed bisimulation, not a quotient}
\begin{defn}[Bisimulation; operational identity $\eq$]
A relation $R:x.V\to y.V\to\mathrm{Prop}$ is a bisimulation between $x$ and $y$ if related vertices
reveal matching members in both directions: for all $a,b$ with $R\,a\,b$,
\[
(\forall a',\ x.E\,a'\,a \Rightarrow \exists b',\ y.E\,b'\,b \wedge R\,a'\,b')\ \wedge\
(\forall b',\ y.E\,b'\,b \Rightarrow \exists a',\ x.E\,a'\,a \wedge R\,a'\,b').
\]
We write $x\eq y$ ($x$ and $y$ are operationally identical) when some bisimulation relates their
points: $\exists R,\ R\,x.\mathrm{pt}\,y.\mathrm{pt}\ \wedge\ R\text{ is a bisimulation}$.
\end{defn}
\begin{prop}[$\eq$ is an equivalence; \texttt{Equiv.refl/symm/trans}]
$\eq$ is reflexive, symmetric and transitive. Each is axiom-free in Lean (no \texttt{propext}, no
\texttt{Quot.sound}, no \texttt{Classical.choice}).
\end{prop}
The methodological point is decisive. To turn ``bisimilar'' into a type-level equality --- to form
the quotient --- one must Skolemise the existential matches into functions; that step is exactly VR's
transit asymmetry (operational$\to$formal is free, formal$\to$operational has no mechanism), and it
is precisely where mathlib's coinductive set construction invokes \texttt{Classical.choice}. We
therefore do not quotient. Identity remains a relation carried with its witness: a performed
bisimulation, not a given completion. This is ``doing, not being'' made into the construction, and
it is why the whole development stays below the choice floor.

\subsection{Membership and strong extensionality}
\begin{defn}[Membership; \texttt{Mem}]
$z\in x$ iff some member-vertex of $x$ is operationally identical to $z$:
$\exists a,\ x.E\,a\,x.\mathrm{pt}\ \wedge\ z\eq x.\mathrm{child}\,a$.
\end{defn}
\begin{thm}[Membership respects identity; \texttt{mem\_congr}]
If $y\eq y'$ then $z\in y\Rightarrow z\in y'$. (Axiom-free.)
\end{thm}
\begin{thm}[Strong extensionality; \texttt{ext}, \texttt{equiv\_iff\_same\_mem}]
Two operational sets are identical iff they have the same members:
$x\eq y \iff \forall z,\ (z\in x \leftrightarrow z\in y)$. Both directions are axiom-free.
\end{thm}
Theorem~2.6 is the heart of the carrier. The forward direction is Theorem~2.5 applied both ways;
the backward direction is the coinductive content --- one builds a bisimulation out of the
member-matching (relate the two roots; relate any further pair whose rooted subtrees are already
identical). It is the AFA-style strong extensionality, kept choice-free precisely because identity
was never quotiented. It is what licenses $\eq$ as a faithful equality of sets.

%======================================================================
\section{The operations}
All operations are instances of one constructor.
\begin{defn}[Family supremum; \texttt{sup}]
For a family $c:I\to\OpSet$, the set $\sup c$ whose members are exactly the $c\,i$ is the pointed
graph with a fresh root revealing each component's point, components kept disjoint.
\end{defn}
\begin{thm}[\texttt{mem\_sup}]
$z\in\sup c \iff \exists i,\ z\eq c\,i$. (Axiom-free.)
\end{thm}
From this single law the standard operations follow; each is axiom-free or depends only on
\texttt{propext}:
\begin{itemize}
\item Empty set $\varnothing$ (\texttt{emptySup}): the empty family. $z\notin\varnothing$
(\texttt{not\_mem\_emptySup}). $\varnothing$ is the nullary operation --- the functionality that
reveals nothing --- matching the base of the VR formal system.
\item Singleton $\{a\}$: $z\in\{a\}\iff z\eq a$.
\item Unordered pair $\{a,b\}$: $z\in\{a,b\}\iff z\eq a\vee z\eq b$.
\item Union $\bigcup x$: $z\in\bigcup x\iff\exists y,\ y\in x\wedge z\in y$ (\texttt{mem\_union}).
\item Binary union $a\cup b$ and successor $\mathrm{succ}\,a:=a\cup\{a\}$:
$z\in\mathrm{succ}\,a\iff z\in a\vee z\eq a$ (\texttt{mem\_succ}).
\end{itemize}
\begin{prop}[Operations respect identity; \texttt{*\_congr}]
\texttt{singleton}, \texttt{pair}, $\cup$ and \texttt{succ} all respect $\eq$ (through Theorem~2.6).
\end{prop}

\subsection{Infinity}
\begin{defn}[Von Neumann naturals and $\omega$; \texttt{vn}, \texttt{omega}]
$\mathrm{vn}\,0:=\varnothing$, $\mathrm{vn}(n{+}1):=\mathrm{succ}(\mathrm{vn}\,n)$;
$\omega:=\sup\mathrm{vn}$.
\end{defn}
\begin{thm}[Infinity; \texttt{mem\_omega}, \texttt{empty\_mem\_omega}, \texttt{omega\_succ\_closed}]
The members of $\omega$ are exactly the $\mathrm{vn}\,n$ up to $\eq$; $\varnothing\in\omega$; and
$\omega$ is closed under successor for any member. (Axiom-free, resp.\ \texttt{propext}.)
\end{thm}

\subsection{Separation and Replacement}
A genuine property or function of sets cannot depend on the particular graph presenting a set; it
must respect $\eq$. With that --- and no decidability assumption --- the two schemas hold.
\begin{thm}[Separation; \texttt{sep}, \texttt{mem\_sep}]
For a species $p$ respecting $\eq$, the set $\{z\in x\mid p\,z\}$ exists and
$z\in\{z\in x\mid p\,z\}\iff z\in x\wedge p\,z$. (\texttt{propext}.)
\end{thm}
\begin{thm}[Replacement; \texttt{repl}, \texttt{mem\_repl}]
For a rule $F$ respecting $\eq$, the image $\{F\,z\mid z\in x\}$ exists and
$w\in\{F\,z\mid z\in x\}\iff\exists z,\ z\in x\wedge w\eq F\,z$. (\texttt{propext}.)
\end{thm}

\subsection{Power set}
A subset of $x$ is selected by a predicate on $x$'s member-vertices; the power set ranges over all of
them. (This is the one operation whose construction reaches into the ambient type theory; see
Remark~3.10.)
\begin{defn}[Subset; power set; \texttt{Subset}, \texttt{subsetOf}, \texttt{powerset}]
$z\subseteq x:\iff\forall w,\ w\in z\Rightarrow w\in x$. For a vertex-predicate $s:x.V\to\mathrm{Prop}$,
$\texttt{subsetOf}\,x\,s$ is the $\sup$ of the member-vertices of $x$ kept by $s$. The power set
$\wp\,x:=\sup_{s:x.V\to\mathrm{Prop}}\texttt{subsetOf}\,x\,s$ ranges over all such predicates.
\end{defn}
\begin{thm}[Power set; \texttt{mem\_powerset}]
$z\in\wp\,x \iff z\subseteq x$. (\texttt{propext}; choice-free.) The backward direction selects the
predicate ``$a$ is a member-vertex whose member lies in $z$'' and closes by strong extensionality
(Theorem~2.6).
\end{thm}
With the power set, \OpSet{} verifies every axiom of ZF without Foundation --- Extensionality,
$\varnothing$, Pairing, Union, Power, Separation, Replacement, Infinity --- machine-checked below the
choice floor; Foundation is the predicate \texttt{IsGrounded} of \S4.
\begin{rem}[The obstruction to power set is impredicativity, not choice]
The completed power set $\wp\,x$ is constructed \emph{choice-free} (Theorem~3.9); in particular
$\wp\,\omega$ exists as a set with $z\in\wp\,\omega\iff z\subseteq\omega$. What this costs is \emph{not}
\texttt{Classical.choice} --- the axiom audit reports \texttt{propext} only --- but
\emph{impredicativity}: the index $x.V\to\mathrm{Prop}$ is the full power of the representative's
vertex type, a move legitimate in Lean's impredicative \texttt{Prop} but outside a predicative
(Brouwer / Martin-L\"of) discipline, and one that \texttt{\#print axioms} cannot detect. So the
operational register's relation to power set is sharper than the ``absent'' of Version~1.0.4's \S7: the
\emph{predicative} content of $\wp\,\omega$ is the becoming --- the non-enumerable diagonal of \S6.1
--- while the \emph{impredicative completed} $\wp\,\omega$ also exists, choice-free, at the cost named
here. The obstruction is isolated: it is impredicativity, not choice.
\end{rem}

%======================================================================
\section{Foundation is a predicate, not a choice}
Classical set theory must decide, at the level of its axioms, whether $\in$-descent always
terminates: ZFC posits Foundation, Aczel's theory posits the Anti-Foundation Axiom (AFA). These are
two incompatible commitments, and a classical theory is forced to pick one. VR is not. On the
operational carrier, well-foundedness is a property of a set's graph, observed and not postulated.
\begin{defn}[Groundedness; \texttt{IsGrounded}]
$x$ is grounded iff its point is accessible under the reveal relation: $\mathrm{Acc}\,(E)\,
x.\mathrm{pt}$ --- every membership descent from $x$ terminates.
\end{defn}
\begin{thm}[Groundedness is a property of the set; \texttt{isGrounded\_iff}]
If $x\eq y$ then $x$ grounded $\iff y$ grounded. (Axiom-free.)
\end{thm}
\begin{thm}[The ZFC fragment is hereditary; \texttt{mem\_grounded}]
A member of a grounded set is grounded. (Axiom-free.)
\end{thm}
Theorem~4.2 makes ``well-founded'' a well-defined predicate on operational sets (transported across
a bisimulation; the key lemma \texttt{acc\_bisim} carries accessibility along $R$). Hence the ZFC
fragment is $\{x\mid x\text{ grounded}\}$ ($\varnothing$ is grounded; the fragment is closed downward
under $\in$), and the ZFA fragment is the whole universe. The Quine atom $A=\{A\}$ is a perfectly
good operational set --- one vertex with a self-loop (\texttt{quine}), satisfying $A\in A$
(\texttt{quine\_self\_mem}) --- and it is not grounded (\texttt{quine\_not\_grounded}); no
anti-foundation axiom is invoked, the graph simply has a cycle. VR thus holds ZFC and ZFA natively,
as fragments of one universe distinguished by a predicate, with uniform identity (witnessed
bisimulation) throughout. The grounded fragment enjoys $\in$-acyclicity: no von Neumann natural is a
member of itself (\texttt{vn\_not\_self\_mem}, by strong induction).

\subsection{Anti-Foundation as a theorem}
Because membership is the edge relation and a member is the same graph re-pointed
(\texttt{child}), \emph{nesting is native}: a set contains a set by an edge, never by copying a set
into another, and self-membership is a self-loop. Aczel's decoration of a graph is therefore almost
definitional here.
\begin{defn}[Decoration; \texttt{decorate}, \texttt{IsDecoration}]
A function $d:V\to\OpSet$ decorates a graph $(V,E)$ if the members of $d\,v$ are exactly the
$d\,w$ with $E\,w\,v$. The canonical decoration is $\texttt{decorate}\,E\,v:=(V,E,v)$ --- the graph
re-pointed at $v$.
\end{defn}
\begin{thm}[Anti-Foundation; \texttt{afa}, \texttt{decorate\_isDecoration}, \texttt{decoration\_unique}]
Every graph has a decoration, unique up to operational identity: $\texttt{decorate}\,E$ decorates
$(V,E)$ --- by \texttt{rfl}, since membership \emph{is} the decoration condition --- and any
decoration $d$ is pointwise $\eq$ to it: for each $v$, the choice-free \emph{exhibited} bisimulation
$R_v\,a\,w:=(d\,v).\texttt{child}\,a\eq d\,w$ witnesses $d\,v\eq\texttt{decorate}\,E\,v$.
\textbf{Fully axiom-free} (\texttt{[]}: no \texttt{propext}, no \texttt{Quot.sound}, no
\texttt{Classical.choice}).
\end{thm}
AFA is thus a theorem, not an axiom; and ZFA is not even a separate commitment --- the Quine atom
(one vertex, a self-loop) and every cycle are simply admissible graphs. With this and
Foundation-as-predicate, \OpSet{} holds ZFC and ZFA natively on one carrier.

\paragraph{The cost of the quotient --- bridge to VR-Sets-ZFA.} The companion work VR-Sets-ZFA
obtains the same anti-foundational content on mathlib's coinductive M-types, by \emph{quotienting} an
extensional bisimulation. There AFA carries the full \texttt{[propext, Classical.choice, Quot.sound]}:
forming the quotient Skolemises the bisimulation's existential witnesses (VR's transit asymmetry,
\S2.2), and uniqueness must select representatives by \texttt{Classical.choice} and run
\texttt{PFunctor.M.bisim} through the choice-pulling destructor \texttt{dest}. The contrast is exact
and is itself a result: \emph{the relational identity is axiom-free precisely where its quotient sits
at the full ceiling} --- the price of turning a performed bisimulation into a given equality is
\texttt{Classical.choice}. We read the two works together as the equivalence
$\mathrm{OSetZFA}\simeq\OpSet/\!\eq$ --- the version here being the choice-free floor of which
VR-Sets-ZFA is the quotient. \textbf{This last equivalence is a prose-level reading, not yet
machine-checked} (in the modality of Conjecture~IV.2): formalising it --- the \texttt{M.corec}
unfolding of a graph, its descent through the quotient, and bijectivity on the $\eq$-classes --- is a
single \texttt{Bridge} file, a candidate for VR-Sets-ZFA~v1.1. Everything else in this paper carries a
named Lean object and its profile.

%======================================================================
\section{Numbers}
The natural numbers are the von Neumann naturals of \S3, genuinely distinct as operational sets.
\begin{thm}[The operational naturals are distinct; \texttt{vn\_inj}]
$\mathrm{vn}\,n\eq\mathrm{vn}\,m\Rightarrow n=m$. (\texttt{propext}.)
\end{thm}
The proof is the classical one read operationally: the members of $\mathrm{vn}\,m$ are exactly
$\mathrm{vn}\,0,\dots,\mathrm{vn}(m{-}1)$, no $\mathrm{vn}\,n$ is a member of itself, and these force
injectivity by the trichotomy of $\N$.

\begin{rem}[Three reals --- stated honestly]
The reals of the VR cycle come in distinct guises that must not be conflated.
\begin{itemize}
\item \textbf{VR-Numbers' $\R_{VR}$} is built from Cauchy sequences with a \emph{computable modulus
of convergence}; it is isomorphic to the field of \emph{computable} reals and is therefore
\emph{countable}. (Its Lean realisation additionally proves an isomorphism with the full classical
line, but that is an artifact of Lean's inability to express computability in a single type, not the
content of the construction --- see VR-Numbers \S VIII.6.) This is the \emph{done}/describable side.
\item \textbf{The operational becoming-continuum} --- the space of lawless choice sequences (and the
Bishop reals over $\Z$ on the Brouwer path) --- is built below the choice floor and is
\emph{genuinely non-enumerable}, choice-free. The non-enumerability is real, not classical: it is the
same phenomenon this paper proves for the set universe (Theorem~6.1, \texttt{universe\_not\_enumerable})
and that the continuum work proves for $\N\to\mathrm{Bool}$ (\texttt{branches\_not\_enumerable}).
\item \textbf{The classical line $\R$} (mathlib) is uncountable in the usual, choice-laden sense.
\end{itemize}
These coexist. The countable computable reals are the performed/describable register; the
non-enumerable becoming-continuum is the operational ``$\R$'' of the Brouwer path; the classical
line is the formal-register referent. Which is canonical for the VR continuum is the programme's open
Bishop/Brouwer question; this work commits the set universe to the Brouwer (non-enumerable) reading.
\end{rem}

%======================================================================
\section{Cardinality, the diagonal, and choice --- the corrected core}
This section replaces the central claims of the earlier \S VI.

\subsection{The operational universe is non-enumerable}
\begin{thm}[Non-enumerability of the universe; \texttt{universe\_not\_enumerable}]
No $\N$-indexed enumeration of operational sets is surjective up to operational identity:
$\neg\,\exists e:\N\to\OpSet,\ \forall x,\ \exists n,\ e\,n\eq x$. The proof is choice-free
(\texttt{propext} only).
\end{thm}
The proof is a self-contained Cantor diagonal inside the operational universe. Given a candidate $e$,
form the diagonal set $D:=\{\,\mathrm{vn}\,n\mid \mathrm{vn}\,n\notin e\,n\,\}$, a legitimate
operational set (a $\sup$ over the subtype of indices $n$ with $\mathrm{vn}\,n\notin e\,n$; no
decidability is needed). By the distinctness of the naturals (Theorem~5.1),
$\mathrm{vn}\,k\in D\iff\mathrm{vn}\,k\notin e\,k$. If $e$ were surjective, some $e\,m\eq D$, whence
by extensionality $\mathrm{vn}\,m\in e\,m\iff\mathrm{vn}\,m\notin e\,m$ --- a contradiction.

This is the Brouwerian content of VR-Sets. The diagonal does not fail and the universe is not
countable. (To be precise: the non-enumerability \emph{theorem} is not itself Brouwerian --- the same
diagonal works classically; what is Brouwerian is its \emph{profile} --- \texttt{propext}-only,
choice-free --- and the contrast with the countable \texttt{Desc} register of \S6.2.) Its meaning is the cut between the \emph{done} and the \emph{becoming}: the countable
register of what has been performed never exhausts the universe of sets as becoming. (In the
operational-continuum work the same phenomenon appears for the power set of $\omega$: the
characteristic functions $\N\to\mathrm{Bool}$ are non-enumerable, choice-free.)

\subsection{The describable register is countable}
\begin{defn}[Finite descriptions; \texttt{Desc}, \texttt{IsDescribable}]
Let \texttt{Desc} be the finite inductive syntax with constructors $\varnothing$, \texttt{pair},
\texttt{union}, $\omega$, and let \texttt{eval}$:\texttt{Desc}\to\OpSet$ be its interpretation. A set
$x$ is describable if $x\eq\texttt{eval}\,d$ for some $d:\texttt{Desc}$.
\end{defn}
\begin{thm}[The describable register is countable; \texttt{describable\_countable}]
There is an $\N$-indexed enumeration reaching every describable set (up to $\eq$), exhibited as
$n\mapsto\texttt{eval}(\texttt{decode}\,n)$.
\end{thm}
\begin{rem}[\texttt{Desc} is a minimal representative of ``finite rule'']
\texttt{Desc} fixes one concrete finite syntax ($\{\varnothing,\texttt{pair},\texttt{union},
\omega\}$, without Separation), so it is deliberately \emph{narrower} than the intuitive ``set given
by a finite rule''. This is exactly why Theorems~6.1 and~6.3 do not collide: the diagonal $D$ of
\S6.1 is itself given by a finite rule, yet it is \emph{not} expressible in \texttt{Desc}. The
contrast is robust under this choice: \emph{any} fixed finite syntax yields a countable register, and
the same diagonal escapes it. \texttt{Desc} is a minimal witness of the schema ``done $=$ countable'',
not an exhaustion of the notion of description.
\end{rem}
\begin{rem}[Honest axiom note]
Theorem~6.3 is the only object in the development that depends on \texttt{Classical.choice}, and the
dependence is \emph{borrowed}, not intrinsic. It enters solely through mathlib's proof that the
natural-number pairing is invertible (\texttt{Nat.unpair\_pair}/\texttt{Nat.pair\_eq\_pair} use
choice; \texttt{Nat.pair} itself is axiom-free). The fact is constructively true; mathlib merely
proves it with incidental choice. We keep the theorem and record the cost. The countability of a
register of descriptions is in any case a statement of the \emph{formal register} --- meta-accounting
over a finite syntax --- where classical reasoning is admitted with its cost made visible.
\end{rem}
Theorems~6.1 and~6.3 together are the corrected \S VI: the earlier ``the operational universe is
countable, hence AC is free'' is false; the correct statement is the contrast --- the describable
(done) register is countable, the universe (becoming) is not.

\subsection{The Axiom of Choice is not a theorem}
The earlier \S VI argued that AC holds because the universe is countable, so every family can be
enumerated and ``the first element of each'' chosen. That argument is unsound: its premise (every
family has a countable enumeration) is exactly what Theorem~6.1 refutes. ``AC is free'' and ``the
continuum is genuinely uncountable'' are two settings of one switch, and the VR cycle has set it on
the Brouwer side. The honest operational picture:
\begin{itemize}
\item \textbf{Dependent choice} is operational and choice-free: a history-dependent rule determines a
branch by recursion (the \texttt{operational\_dependent\_choice} of the continuum work).
\item \textbf{Weak continuity (WC-N)} holds operationally for finite-information functionals
(machine-checked there) and is classically false --- the cleanest two-register exhibit.
\item \textbf{Full AC} --- a selection over an uncountable family with no rule --- has no operational
correlate. It is a label of the formal register, not an operational theorem.
\end{itemize}

\subsection{Why Banach--Tarski does not arise --- through conservativity, not counting}
The earlier text located the absence of Banach--Tarski in the countability of the universe. With
countability gone, the explanation must change --- and the correct one is sharper and ties VR-Sets to
VR-Forms. The Banach--Tarski construction is a form applied to a form ($T\to T$): in the formal
register it builds a completed uncountable totality (the sphere as a set of points) with a selection
by full AC (the Vitali set). Neither ingredient has an operational correlate: the completed
uncountable totality is a formal-register referent, and full AC a formal-register label. By the
invariance of the operational under the application of form --- the ground of VR-Forms'
conservativity result, itself machine-checked --- a formal-register construction does not transit to
the operational register: form has no power to manufacture an operational object where none was
built. The paradox lives entirely in the formal register and never returns through the operational
floor. Its absence is a conservativity phenomenon, not a cardinality accident.

There is, moreover, a precise classical anchor. Solovay (1970) showed that
$\mathrm{ZF}+\mathrm{DC}+$``every set of reals is Lebesgue measurable'' is consistent (relative to an
inaccessible cardinal); in that model the Banach--Tarski paradox \emph{provably fails}. The
operational regime --- dependent choice yes, full AC no (\S6.3) --- lands in exactly that territory.
So the absence is not only a conservativity phenomenon but one with known metamathematical backing:
Banach--Tarski is not derivable in the choice regime the operational register supports. (This is an
anchor, not an identification --- the operational register is not claimed to \emph{be} Solovay's
model; the point is that its choice profile sits where the paradox is known to be underivable.)

%======================================================================
\section{Position and honest boundaries}\label{sec:position}
\paragraph{What is established.} The operational universe of sets is built without classical
scaffolding: a set is a revealing functionality, identity is a witnessed bisimulation (no quotient),
and on this carrier membership, strong extensionality, $\varnothing$, pairing, union, the power set,
$\omega$, Separation and Replacement, the distinctness of the naturals, Anti-Foundation (decoration of
any graph, unique up to $\eq$; axiom-free), and the predicate character of
well-foundedness are all machine-checked --- almost entirely axiom-free or \texttt{propext}-only (the
power set is \texttt{propext}, choice-free, though impredicative; Remark~3.10). \OpSet{} thus verifies
every axiom of ZF without Foundation and locates the ZFC and the ZFA universes as two fragments of one
carrier (the precise, two-register sense in which it holds them is stated next). The
corrected cardinality core (universe non-enumerable; describable register countable) is
machine-checked, the substantive half choice-free.

\paragraph{The two registers, and the exact relation to ZFA.} The relation of \OpSet{} to the
classical theories is stated most precisely through VR's two registers. VR's \emph{operational}
register is an unconditional, choice-free constructive ground; its \emph{formal} register admits --- as
a first-class but conditional, cost-tagged act --- any consistent posit (the axiom of choice, actual
infinity, the uncountable). \textbf{With the formal register on} (classical logic, choice), the set
layer of \OpSet{} \emph{realises} \emph{Aczel's} ZFA (that is, $\mathrm{ZF}^-$ with the
Anti-Foundation Axiom --- non-well-founded sets by self-membership, not by urelements) and is
therefore \textbf{equal to ZFA in strength: equiconsistent, at the level of ZF} --- neither stronger
nor weaker. Within it ZFC is recovered as the well-founded inner class $\{x : \texttt{IsGrounded}\,x\}$.
\textbf{Without the formal register}, the operational core is strictly \emph{weaker} --- constructive
and choice-free (nearer $\mathrm{CZF}+\mathrm{AFA}$), not full classical ZFA. Three provisos keep this
exact. (i) It is realisation, not identity of theories: \OpSet{} is a model that satisfies ZFA's
axioms, so it decides statements (e.g.\ CH) that ZFA leaves open; ``$=\text{ZFA}$'' is at the level of
axioms and consistency strength, not of the deductive closure. (ii) The equality is for the
\emph{set layer}; VR as a whole is a proper enrichment of ZFA --- it adds the operational register and
a machine-checked cost ledger that plain ZFA has not. (iii) The one place VR is strictly \emph{better}
than ZFA, at equal strength, is the price of nesting: Anti-Foundation is here a \emph{theorem},
axiom-free (\S4.1), where ZFA must \emph{posit} it --- the same universe, obtained at lower axiomatic
cost, because the carrier (pointed graphs up to bisimulation) makes self-membership definitional rather
than assumed. (The pluralist framing --- consistent theories as legitimate, ranked by strength --- is
not itself novel, cf.\ Hilbert, the set-theoretic multiverse; what is new is the sharp
operational/formal line, native choice-free Anti-Foundation, and the per-posit machine-checked ledger.)

\paragraph{What is open or deferred.}
\begin{itemize}
\item \emph{Bishop vs Brouwer.} The computable $\R_{VR}$ (countable, done) and the non-enumerable
becoming-continuum coexist (Remark~5.2); which is canonical for the VR continuum is the programme's
open question. This version commits the set universe to the Brouwer reading, consistently with the
operational-continuum work.
\item \emph{The metalanguage.} Lean treats the underlying types classically; the operational reading
is carried in the metalanguage and tracked by axiom profile. ``Below the choice floor'' is a
statement about the construction's audited dependencies, exhibited build-time, not a claim that the
ambient logic is intuitionistic.
\item \emph{The borrowed choice.} The countability of the describable register inherits one
\texttt{Classical.choice} from mathlib's pairing inversion (Remark~6.5); removing it would mean
re-deriving the arithmetic choice-free, orthogonal to the foundations and left aside.
\item \emph{Power set --- impredicative, not choice-laden.} A completed power set $\wp\,x$ \emph{is}
constructible choice-free (\S3.3, Theorem~3.9), but only impredicatively (Remark~3.10): the obstruction
is impredicativity, not choice. The \emph{predicative} content of $\wp(\omega)$ remains the becoming
--- the non-enumerable diagonal of \S6. CH is a statement about a formal-register referent; it is read
philosophically, not coded.
\item \emph{Impredicativity is invisible to the axiom audit (P-1).} Remark~3.10 isolates a new kind of
boundary: the power set is choice-free yet impredicative, and \texttt{\#print axioms} cannot see the
impredicativity --- it reports \texttt{propext} only. The operational register therefore needs a
\emph{second dimension} of accounting beyond the axiom profile; naming it (P-1) is the honest cost of
the power set, alongside B-1/B-2 (the choice floor in $\R$ / \texttt{Fin}) and TR-R1 (the cost of the
quotient).
\item \emph{The setoid tax.} Declining the quotient is a permanent discipline, not a one-time saving:
every operation carries a congruence lemma (\texttt{*\_congr}, \texttt{mem\_congr},
\texttt{isGrounded\_iff}), and every downstream use of \OpSet{} pays it. This is the named price of
staying below the choice floor --- relational identity in exchange for a standing congruence
obligation. We pay it deliberately, in keeping with the cycle's discipline of naming prices.
\end{itemize}
\paragraph{Relation to the earlier version.} Versions through 1.0.2 stand as the classical-surrogate
presentation (the ZFC axioms realised over \texttt{ZFSet}); their \S VI cardinality claims are
superseded by \S6 here. The mathematics genuinely machine-checked there (the ZFC axioms as closure
facts over the surrogate) is unaffected; what is corrected is the operational reading of cardinality
and choice.

%======================================================================
\section*{Acknowledgement of AI assistance}
\addcontentsline{toc}{section}{Acknowledgement of AI assistance}
This preprint and its companion Lean~4 formalisation were prepared with the assistance of Claude
(Anthropic), in the Variant~A architecture used across the VR cycle: a human curator directing a
model that served as architect and implementer. All mathematical content, definitions, theorems, and
foundational positions are due to the author, Vitaly Reznik. The operational set universe
(\texttt{VRCycle.SetsOp}), the correction of \S VI, the operational power set (\S3.3), and this
Version~1.1.0 were developed with Claude Opus~4.8.

\section*{References}
\addcontentsline{toc}{section}{References}
\begin{itemize}
\item P. Aczel, \emph{Non-Well-Founded Sets}, CSLI Lecture Notes 14, 1988.
\item E. Bishop, \emph{Foundations of Constructive Analysis}, McGraw--Hill, 1967.
\item L. E. J. Brouwer, collected works on choice sequences and the continuum.
\item R. M. Solovay, \emph{A model of set-theory in which every set of reals is Lebesgue measurable},
Annals of Mathematics \textbf{92} (1970), 1--56.
\item J. von Neumann, \emph{Zur Einf\"uhrung der transfiniten Zahlen}, 1923.
\item V. Reznik, \emph{VR. A Formal System: An Axiom-Free Reconstruction of Arithmetic from Operations Alone.} Zenodo, 2026. DOI 10.5281/zenodo.20633314 (concept 10.5281/zenodo.20212091); Lean companion DOI 10.5281/zenodo.20324240.
\item V. Reznik, \emph{VR-Numbers: Operational Superstructures \dots over the Natural Numbers of VR.} Zenodo, 2026. DOI 10.5281/zenodo.20633626 (concept 10.5281/zenodo.20272025).
\item V. Reznik, \emph{VR-Forms: An Operational Two-Register Apparatus over VR-Sets.} Zenodo, 2026. DOI 10.5281/zenodo.20355939.
\item V. Reznik, \emph{VR-Sets-ZFA.} Zenodo, 2026. DOI 10.5281/zenodo.20369346.
\item The VR Cycle Lean~4 formalisation, \url{https://github.com/inventor1975/VRCycle}, module \texttt{VRCycle.SetsOp}.
\end{itemize}

\end{document}
